\chapter{Electrogravitational Bilinear Form and Cross-Response Equations}
\label{chap:forma-electrogravitatoria}

\section{Internal source coordinates}

Let \(Z\) be the impedance section of the vacuum and let \((Q,\ell)\) be the pair
from the preceding chapter. To a charge \(q\) and an energy \(E\) we associate

\[
\xi_q(q)=\sqrt{\frac{Z}{4\pi\hbar}}\,q,
\qquad
\xi_E(E)=\frac{E}{Q}.
\]

The internal source vector is

\[
X(q,E)=
\begin{pmatrix}\xi_q(q)\\ \xi_E(E)\end{pmatrix},
\qquad
\eta_{1,1}=
\begin{pmatrix}1&0\\0&-1\end{pmatrix}.
\]

The form \(X_1^{\mathsf T}\eta_{1,1}X_2\) separates electrical orientation and
gravitational closure within a single dimensionless plane.

\section{Static unification of the potentials}

\begin{theorem}[Electrogravitational form]
\label{thm:rm-emg-potential}
If \(Zc=\varepsilon^{-1}\), \(E_i=m_ic^2\) and
\(Q\ell=\hbar c\), then
\begin{equation}
\boxed{
V_{12}(r)=\frac{Q\ell}{r}
X(q_1,E_1)^{\mathsf T}\eta_{1,1}X(q_2,E_2)}
\label{eq:rm-emg-potential}
\end{equation}
is exactly the sum of the Coulomb potential and the Newton potential.
\end{theorem}

\begin{proof}
The electric term is
\[
\frac{Q\ell}{r}\frac{Zq_1q_2}{4\pi\hbar}
=\frac{Zcq_1q_2}{4\pi r}
=\frac{q_1q_2}{4\pi\varepsilon r}.
\]
The term gravitational is
\[
-\frac{Q\ell}{r}\frac{E_1E_2}{Q^2}
=-\frac{\ell}{Q}\frac{E_1E_2}{r}
=-\frac{G}{c^4}\frac{m_1m_2c^4}{r}
=-\frac{Gm_1m_2}{r}.
\]
\end{proof}

The associated radial force takes the form

\begin{equation}
F_{12}(r)=\frac{Q}{\ell}
\left(\frac{\ell}{r}\right)^2
X_1^{\mathsf T}\eta_{1,1}X_2.
\label{eq:rm-emg-force}
\end{equation}

The factor \(Q/\ell=\mathscr T_P\) is the universal tension; the
dimensionless form determines the intensity and sign of its publication.

\section{Cone of extremality of constitutive responses and classification of its rays}

For a single source, the discriminant

\[
\Delta_{\rm EMG}(q,E)=E^2-\mathcal Q_e(q)^2,
\qquad
\mathcal Q_e(q)=Q\sqrt{\frac{Z}{4\pi\hbar}}\,q,
\]

induce the three sectors

\[
\Delta_{\rm EMG}>0,
\qquad
\Delta_{\rm EMG}=0,
\qquad
\Delta_{\rm EMG}<0.
\]

The condition null is equivalent to

\begin{equation}
E^2=\mathcal Q_e(q)^2
\quad\Longleftrightarrow\quad
Gm^2=\frac{q^2}{4\pi\varepsilon}.
\label{eq:rm-rn-extremality}
\end{equation}

It is the condition of electrogravitational extremality in the Reissner--Nordström
normalization. Here the trit acquires a static realization:
subextremality, extreme threshold and superextremality are the three signs of the
same quadratic form.

For the elementary charge constructed by

\[
e^2=\frac{4\pi\alpha_{\rm HMT}\hbar}{Z},
\]

the coordinate satisfies

\begin{equation}
\boxed{\xi_q(e)^2=\alpha_{\rm HMT}.}
\label{eq:rm-electron-electric-coordinate}
\end{equation}

This equality locates the central charge in the source plane without using the
electron mass or importing its complete theory.

\section{Einstein--Maxwell Coefficient}

Write

\[
\kappa=\frac{8\pi G}{c^4}=8\pi\mathscr C_P.
\]

In a convention where the electromagnetic tensor is expressed as
\(T_{\mu\nu}^{\rm EM}=\mu^{-1}\mathcal Q_{\mu\nu}(F,g)\), the gravitational coefficient
of the quadratic form is \(\kappa/\mu\). Using

\[
\ell^2=\frac{G\hbar}{c^3},
\qquad
Z=\mu c,
\qquad
e^2=\frac{4\pi\alpha\hbar}{Z},
\]

we obtain

\begin{equation}
\boxed{
\frac{\hbar^2}{8\pi\ell^2}
\frac{(\kappa/\mu)}{e^2}
=\frac{1}{4\pi\alpha_{\rm HMT}}.}
\label{eq:rm-einstein-maxwell-alpha}
\end{equation}

\begin{proof}
First,
\[
\frac{\kappa}{\mu}
=\frac{8\pi G}{c^4\mu}
=\frac{8\pi\ell^2}{\hbar\mu c}
=\frac{8\pi\ell^2}{\hbar Z}.
\]
Al dividir by \(e^2=4\pi\alpha\hbar/Z\) and multiplicar by
\(\hbar^2/(8\pi\ell^2)\) remains \(1/(4\pi\alpha)\).
\end{proof}

At the rationalized electroweak interface,

\[
g^2=\frac{4\pi\alpha}{1-D_A},
\qquad
g'^2=\frac{4\pi\alpha}{D_A},
\]

and therefore

\begin{equation}
\boxed{
\frac{\hbar^2}{8\pi\ell^2}
\frac{(\kappa/\mu)}{e^2}
=\frac1{g^2}+\frac1{g'^2}.}
\label{eq:rm-cross-einstein-ew}
\end{equation}

The equality requires the gravity, vacuum, charge, and gauge
mixing readers to commute. Within the present graph, \(e,g,g'\) descend
from \(\alpha_{\rm HMT}\); the equation is a cross-check, not a second
derivation independent of \(\alpha\).

\section{Polygon of fine structure publications}

The dodecaphasic macro-object is published through the system

\begin{equation}
\boxed{
\alpha_{\rm HMT}
=\frac{Ze^2}{4\pi\hbar}
=\frac{ZG_0}{4}
=\frac{Z}{2R_K}
=-\frac{9}{100\pi}\log D_A
=\frac{(1-D_A)g^2}{4\pi}
=\frac{D_Ag'^2}{4\pi}.}
\label{eq:rm-alpha-polygon}
\end{equation}

The first causal construction of \(\alpha\) remains the signed wheel,
the seal, and the carry. The subsequent equalities are metrological publications
or compatibility checks. Elimination of variables yields

\[
(1-D_A)g^2=D_Ag'^2
=\frac{Ze^2}{\hbar}
=\pi G_0Z
=\frac{2\pi Z}{R_K}.
\]

The dual form is

\[
\frac1{g^2}+\frac1{g'^2}
=\frac{\hbar}{Ze^2}
=\frac{R_K}{2\pi Z}
=\frac1{\pi G_0Z}.
\]

Each metrological cycle must be marked as a \emph{round trip}: if \(e\),
\(R_K\), or \(G_0\) were defined from \(\alpha\), their inverse substitution
checks coherence but adds no selective information.

\section{Covariance of action inversion}

Under the bidirectional transformation,

\[
\frac{e_{\rm pre}^2}{e_{\rm ret}^2}=R_{\rm act},
\qquad
\frac{G_{\rm pre}}{G_{\rm ret}}=R_{\rm act}^{-1}.
\]

Therefore,

\begin{equation}
\boxed{G_{\sigma}e_{\sigma}^2=\text{sheet constant}.}
\label{eq:rm-ge2-invariant}
\end{equation}

In constitutive form,

\[
\frac{Ze_{\sigma}^2G_{\sigma}}
{c^5t_0^2}
=4\pi\theta_{108}\alpha_{\rm HMT},
\qquad
\theta_{108}=\left(\frac{54}{\pi}\right)^2.
\]

The charge preserves the orientation of action, and the gravitational response
compensates for it; their product belongs to the shared geometry.

\Needspace{12\baselineskip}
\section{Tensorial representation of the gravitational carrier and transverse reduction of its degrees of freedom}

The bilinear form \eqref{eq:rm-emg-potential} closes the static sector. Its
promotion to complete dynamics requires a realizer

\[
\mathcal R_{\rm EM}:
\mathcal X_{\TPK}^{\rm enr}
\longrightarrow
(e^I{}_{\mu},\omega^{IJ}{}_{\mu},A_\mu)
\]

that preserves the Bianchi identity, Maxwell's equations,
current conservation, and the boundary terms. The identity
\eqref{eq:rm-cross-einstein-ew} fixes the coefficient that this realizer must
publish; it does not replace its differential equations.

\begin{falsifierbox}
The electrogravitational form is refuted by: a different coefficient in Coulomb or
Newton upon expanding \eqref{eq:rm-emg-potential}; a null cone that fails to reproduce
\eqref{eq:rm-rn-extremality}; the loss of
\eqref{eq:rm-einstein-maxwell-alpha}; or the use of a publication
descended from \(\alpha\) as a retroactive selector of its dodecaphase branch.
\end{falsifierbox}

\paragraph{Causal structure of the electrogravitational confluence.}
The composition proved in this chapter is organized into five typed stages. The enriched state first supplies the phase \(\alpha_{\rm HMT}\) and its orientation. The constitutive operator then publishes \(e\), \(g\), and \(g'\) in their respective codomains. The gravitational closure determines \(G\) and the chronology \(\theta_{108}\). Identities \eqref{eq:rm-emg-potential}--\eqref{eq:rm-cross-einstein-ew} coordinate these publications through a common coefficient. Finally, \(\mathcal R_{\rm EM}\) transports the result to the fields \((e^I{}_{\mu},\omega^{IJ}{}_{\mu},A_\mu)\), preserving the differential identities and boundary terms. This factorization fixes the antecedence of the HMT construction with respect to its tensorial realization and prepares the constitutive relations developed in the following chapter.
